\section{GPU Acceleration}
\label{sec:gpu}

Multiway evolution is an irregular workload for a GPU. The frontier's width depends on the data
and is not known until the step that produces it, and every new state must be canonicalized before
it can be compared with others, by the same IR code the host runs, whose control flow depends on
the data. The kernel spreads IR's loops across the lanes of a warp where the loops allow it. A
kernel launch per step would put a host round trip between every stage, so the engine runs one
resident kernel that keeps the frontier on the device for the whole evolution.
Appendix~\ref{app:gpu-impl} specifies the device canonicalization, matching, work queues,
termination and memory management.

\subsection{Resident Kernel}

A kernel launch has a fixed cost, which dominates the early steps of an evolution, where the
frontier is small. The engine runs a single persistent kernel for the whole evolution, with the
frontier, the match queues and the epoch counter kept in device memory, so a step costs a
device-side barrier rather than a launch. The match, rewrite and state queues are resident, and a
new state becomes a match task without returning to the host. The resident kernel is
$\PersistentSpeedupFour\times$ faster than a per-step \texttt{evolve()} at depth four, and the two
take the same time from depth \PersistentTieDepthWord{} on, where the evolution itself outweighs
the launch cost (Section~\ref{sec:benchmarks}).

One block matches one (state, rule) pair, and its threads divide the candidates for the first
pattern edge. Later pattern edges take their candidates from a vertex index through a variable
that is already bound, as described in Appendix~\ref{app:gpu-impl}.

\subsection{Design Constraints}
\label{sec:gpu-constraint}

Nothing crosses the host--device boundary while the kernel runs. Rules and the initial state are
uploaded and the queues seeded before the launch, and results are read after it; in between, the
device decides what work exists, which worker takes it, and when the evolution has finished. Any
round trip during the run would cost as much as the per-step round trip the design removes.

This has two consequences.

\emph{The kernel cannot grow a pool.} Growing a pool needs a host allocation. When a pool
overflows, the kernel records the overflow and returns the work completed so far; the host then
re-runs the evolution with doubled capacities, up to eight times, and returns the last partial
result with a warning if it still overflows (Section~\ref{sec:capacity}).

\emph{Canonicalizer scratch cannot be sized from a batch.} The resident kernel has no batch: states
arrive continuously and the largest is not known until the run ends. Scratch pre-sized from a
configured maximum, with a coarser hash above it, would let two non-isomorphic states share a key
whenever one was too large to canonicalize exactly. Instead each state takes scratch from a
device-side arena, sized from its own vertex, edge and occurrence counts. Exhausting the arena is a
capacity overflow, handled as above; the exact path never falls back to a coarser identity.
